\Needspace{18\baselineskip}
\subsection{Auxiliary compact complex factors}
\label{subsec:general-auxiliary-product-models}

So far, the auxiliary factor that makes the background nonclosed has been a curve. By \cref{lem:auxiliary-slice-envelopes}, only the semipositivity of the auxiliary term and its vanishing on the slices parallel to the curve product matter. The auxiliary factor can therefore be any compact complex manifold of positive dimension, K\"ahler or not, and the potentials and their polar metric still depend only on the curve coordinates and the fixed divisor of weight one.

Keep $Y$, $\sigma_i$, $B_i$ and $\gamma_\sigma$, as well as the data $(Z,\kappa,D,\psi)$ of the divisor of weight one. Let $T$ be a connected compact complex manifold of dimension $\ell\geq1$, and put $B=Y\times T$, with projections $\pi$ to $Y$ and $\operatorname{pr}_T$ to $T$. Choose a smooth semipositive real $(1,1)$-form $\alpha$ on $B$ whose restriction to every slice $Y\times\{t\}$ is zero. By semipositivity, the tangent spaces of these slices lie in the kernel of $\alpha$, so in product coordinates $\alpha$ has only coefficients in the $T$-directions. On $M=B\times Z$, let $\varpi,r$ be the projections, and put $p=\pi\circ\varpi$ and
\[
 \Gamma_\alpha=\pi^*\gamma_\sigma+\alpha,\qquad
 \Theta_\alpha=\varpi^*\Gamma_\alpha+r^*\kappa,
 \qquad \Psi=\psi\circ r.
\]
We assume neither that $\alpha$ is closed nor that $M$ has the bounded mass property. For $T=C_0$ and $\alpha=(f\circ\pi)\operatorname{pr}_T^*\tau$, $\Theta_\alpha$ is the background $\Theta_{\sigma,f}$ of \cref{subsec:arbitrary-curve-product-metrics}.

\Needspace{16\baselineskip}
\begin{lemma}\label{lem:general-auxiliary-product-projections}\index{Envelope!compact complex auxiliary factor}
For $u_i\in\PSH(C_i,\sigma_i)$ put
\[
 u=\sum_i u_i\circ\pi_i,\qquad U=\Psi+u\circ p,
 \qquad q_i=P_{\sigma_i}[u_i],\quad
 r_i=P_{\sigma_i}[u_i]_{\mathcal I}.
\]
Then
\[
 \begin{aligned}
 P_{\Theta_\alpha}[U]&=\Psi+\sum_iq_i\circ\pi_i\circ p,\\
 P_{\Theta_\alpha}[U]_{\mathcal I}
     &=\Psi+\sum_ir_i\circ\pi_i\circ p.
 \end{aligned}
\]
Every ordinary representative of either of these envelopes has zero extended mass, so their upper and lower singularity volumes are zero.
\end{lemma}
\begin{proof}
By \cref{prop:separated-curve-obstacle-projections}, $P_{\gamma_\sigma}[u]=\sum_iq_i\circ\pi_i$, and $P_{\gamma_\sigma}[u]_{\mathcal I}=\sum_ir_i\circ\pi_i$ satisfies the divisorial hypothesis of \cref{lem:auxiliary-slice-envelopes}. That lemma gives
\[
 P_{\Gamma_\alpha}[u\circ\pi]=\sum_iq_i\circ\pi_i\circ\pi,\qquad
 P_{\Gamma_\alpha}[u\circ\pi]_{\mathcal I}=\sum_ir_i\circ\pi_i\circ\pi.
\]
\Cref{prop:fibration-obstacle-reduction}, applied to $M\to B$ and the divisor $B\times D$ with weight one, whose induced zero obstacle is zero because $\sup_Z\psi=0$, gives both identities upstairs, and its mass assertion gives the volume assertions.
\end{proof}

For normalized ordinary curve models $q_i$, write
\[
 \Phi_\alpha(q)=\Psi+\sum_iq_i\circ\pi_i\circ p,
\]
let $\mathcal S_\alpha$ be the locus of these potentials, and give it the metric
\[
 \rho_\alpha(\Phi_\alpha(q),\Phi_\alpha(\widetilde q))
       =\sum_i\|\mu_i-\widetilde\mu_i\|_{\mathrm{TV}},
\]
where $\mu_i,\widetilde\mu_i$ are the polar curvatures of $q_i,\widetilde q_i$.

\Needspace{18\baselineskip}
\begin{proposition}\label{prop:general-auxiliary-product-metric}\index{Metric completeness!compact complex auxiliary factor}\index{Model potential!auxiliary factor independence}
The map $\Phi_\alpha(q)\mapsto\Phi_\sigma(q)$ is an isometry from $\mathcal S_\alpha$ onto the locus $\mathcal S_\sigma$ of \cref{prop:arbitrary-curve-product-metric}, and it maps the ideal subspaces and the models with analytic singularities onto each other. Consequently $\mathcal S_\alpha$ is complete and geodesic, with diameter $2\sum_iB_i$. Its ideal subspace is closed, geodesically convex and equal to the metric closure of the models with analytic singularities. Ideal projection is the unique nearest-point retraction onto this subspace and is $1$-Lipschitz, and the distance to the subspace is the sum of the nonatomic polar masses.

Changing $T$ or $\alpha$ does not change this abstract metric space. All the area comparisons and the contractions commuting with ideal projection of \cref{prop:arbitrary-curve-product-metric} hold, even for different auxiliary data. The upper-volume comparison vanishes on the locus, and metric convergence implies $L^1(M)$ convergence for every smooth positive volume form. Suppose that $T$ and $Z$ are fixed and that the curve backgrounds converge uniformly to backgrounds of zero area. Then, for any choice of the semipositive terms $\alpha$, the potentials converge in $L^1(M)$, uniformly over the model tuples, to the sum of $\Psi$ and the pullbacks of the normalized smooth potentials of the limiting curve backgrounds, as in \cref{cor:arbitrary-product-zero-area-collapse}.
\end{proposition}
\begin{proof}
Applied to model tuples, the lemma shows that every $\Phi_\alpha(q)$ is a normalized ordinary model and computes its ideal projection coordinatewise. Restrictions to curve sections on which the other potentials are finite determine the tuple up to constants, and the normalization of the coordinates determines these constants. Thus the map is well defined and bijective, and the formulas for the polar curvatures make it an isometry preserving ideal projection.

Each potential in either locus is pulled back from $Y\times Z$. Analytic presentations pull back under the projection that forgets the auxiliary factor. Conversely, by \cref{prop:analytic-type-calculus}, the restriction of a potential with analytic singularities to a section with fixed auxiliary coordinate has analytic singularities and is not identically $-\infty$. Hence a potential in either locus has analytic singularities exactly when the common potential on $Y\times Z$ does, and the isometry maps the analytic models onto each other. The statements on the complete metric, the ideal closure, geodesics, nearest points and area comparisons now follow from \cref{prop:arbitrary-curve-product-metric}.

A pairwise maximum is $\Psi$ plus the pullback of the maximum of two $\Gamma_\alpha$-plurisubharmonic functions on $B$. The reduction along the divisor of weight one gives zero upper volume for it and for both models, so the upper-volume comparison vanishes. For $L^1$ convergence, integration along the fixed compact fibers bounds the norm of each pulled-back curve potential by a constant times its norm on the curve. The same estimate and \cref{thm:curve-zero-area-collapse} give the asserted uniform limit at zero area. Neither estimate involves $\alpha$.
\end{proof}

\Needspace{16\baselineskip}
\begin{corollary}\label{cor:general-auxiliary-all-obstacle-criterion}\label{cor:auxiliary-product-all-obstacle-criterion}\index{Envelope!all-obstacle compact auxiliary criterion}\index{Envelope!nonclosed separated product criterion}
For the arbitrary curve potentials $u_i$ of the lemma, put $R=\Psi+\sum_ir_i\circ\pi_i\circ p$. The following are equivalent: $P_{\Theta_\alpha}[U]=P_{\Theta_\alpha}[U]_{\mathcal I}$; for each $i$, the polar part of the curvature measure of $u_i$ equals its atomic part; $U\in\mathcal E_{\mathrm{env}}(M,\Theta_\alpha,R)$; and the ordinary and ideal envelopes of $U$ agree below every smooth real obstacle on $M$.
\end{corollary}
\begin{proof}
By the lemma and the normalization of the coordinates, equality of the envelopes below zero is equivalent to $q_i=r_i$ for every $i$. By \cref{cor:curve-arbitrary-background-criterion}, this is equivalent to equality of the polar and atomic measures, and to $u_i\in\mathcal E_{\mathrm{env}}(C_i,\sigma_i,r_i)$ for every $i$. In that case \cref{lem:product-relative-envelope-classes} gives $u\in\mathcal E_{\mathrm{env}}(Y,\gamma_\sigma,\sum_ir_i\circ\pi_i)$. Since $\alpha\geq0$, pullbacks of stretching competitors are $\Gamma_\alpha$-plurisubharmonic, so $u\circ\pi\in\mathcal E_{\mathrm{env}}(B,\Gamma_\alpha,\sum_ir_i\circ\pi_i\circ\pi)$, and \cref{prop:fibration-relative-envelope-class} gives $U\in\mathcal E_{\mathrm{env}}(M,\Theta_\alpha,R)$. By the lemma, $R=P_{\Theta_\alpha}[U]_{\mathcal I}$, so \cref{prop:ideal-obstacle-reprojection} gives equality below every smooth obstacle. Taking the zero obstacle recovers the first condition, which proves the equivalence.
\end{proof}

\Needspace{16\baselineskip}
\begin{corollary}\label{cor:general-auxiliary-background-boundaries}\index{Model potential!nonbig auxiliary product}
If $\alpha=0$ or some $B_i=0$, the background $\Theta_\alpha$ is not big. If all $B_i>0$ and $\alpha$ is positive on every nonzero vector tangent to the factor $T$, it is big. In particular, for a Hermitian form $\omega_T$ on $T$ and a smooth function $h\geq0$ on $Y$, the choice $\alpha=(h\circ\pi)\operatorname{pr}_T^*\omega_T$ gives a background that is not closed when $h$ is nonconstant, and big when $h>0$ and all $B_i>0$.
\end{corollary}
\begin{proof}
For $\alpha=0$, the background is pulled back from $Y\times Z$, and by \cref{prop:psh-descent-submersion} every potential in it descends to $Y\times Z$. The corresponding currents have no component in the $T$-directions, so they cannot dominate a positive multiple of a Hermitian form on $M$. Thus the background is not big.

If $B_i=0$, choose any Hermitian form $\omega_T$ on $T$ and write $s\colon M\to T$ for the projection. The positive $(m+\ell+k-1,m+\ell+k-1)$-form
\[
 L_i=p^*\left(\bigwedge_{j\ne i}\pi_j^*\eta_j\right)
       \wedge s^*(\omega_T^\ell)\wedge r^*(\kappa^k)
\]
is closed, because $\omega_T^\ell$ has top degree on $T$. Its pairing with $\Theta_\alpha$ is zero, since the only possible contribution contains the integral of $\sigma_i$, whereas its pairing with a positive product Hermitian form is strictly positive. By Stokes' theorem, no current in the background is bounded below by a positive multiple of a Hermitian form.

When all areas are positive, write $\sigma_i=B_i\eta_i/A_i+\ddc\chi_i$. Subtracting the smooth pullback of $\sum_i\chi_i\circ\pi_i$ gives the current $p^*\sum_i(B_i/A_i)\pi_i^*\eta_i+\varpi^*\alpha+r^*\kappa$. Positivity of $\alpha$ in the $T$-directions makes this form strictly positive, which proves bigness. For the stated choice of $\alpha$, the exterior derivative of $\Theta_\alpha$ has the component $\mathrm dh\wedge\omega_T$, with one direction in $Y$ and two in $T$. It is nonzero for nonconstant $h$ and cannot be canceled by the component with three directions in $T$. Hence $\mathrm d\Theta_\alpha\ne0$.
\end{proof}

The proofs use no bound on the ambient masses, so the auxiliary manifold may fail the bounded mass property; an explicit example follows.
